\section{Motywacja}
Motywacją do podjęcia tego tematu jest między innymi globalny wzrost zainteresowania technologiami wykorzystującymi sztuczną inteligencję (ang.~\gls{AI}).
Według raportu przygotowanego przez Uniwersytet Stanforda, generatywna sztuczna inteligencja (ang.~\gls{GenAI}) została zaadaptowana przez 53\% populacji w ciągu trzech lat, co czyni ją jedną z najszybciej wdrażanych technologii w historii~\cite{ai_index_2026}.
Ważną częścią tych systemów stają się duże modele językowe (ang.~\gls{LLM}), które używane są do tworzenia zaawansowanych wyszukiwarek, chatbotów czy narzędzi wspomagających wytwarzanie oprogramowania~\cite{nist_taxonomy_2025}.
Stosowanie modeli językowych na skalę masową sprawia, że zagwarantowanie ich bezpieczeństwa ma wyjątkową wagę.
Spośród zidentyfikowanych zagrożeń dla systemów opartych na LLM, szczególne znaczenie mają ataki typu wstrzykiwanie poleceń (ang. prompt injection). Polegają na celowym manipulowaniu zachowaniem lub odpowiedzią modelu językowego w sposób niezamierzony poprzez polecenie użytkownika~\cite{owasp_llm_top10}.
Istotność tego ryzyka potwierdza klasyfikacja przygotowana przez organizację Open Worldwide Application Security Project (OWASP), która w swoim zestawieniu uznaje tego typu ataki za największe zagrożenie dla bezpieczeństwa we współczesnych aplikacjach wykorzystujących modele językowe~\cite{owasp_llm_top10}.
Z tego względu główną motywacją do podjęcia niniejszych badań jest chęć analizy i porównania skuteczności metod mitygacji tej krytycznej podatności w obliczu szybkiego tempa integracji systemów GenAI i potrzeby zapewnienia wysokiego poziomu bezpieczeństwa dużych modeli językowych.